\documentclass[10pt,twocolumn,letterpaper]{article}

% ---------------------------------------------------------------
% Include basic CVPR package in rebuttal mode

\usepackage[rebuttal]{cvpr}

% ---------------------------------------------------------------
% Other packages

% Commonly used abbreviations (\eg, \ie, \etc, \cf, \etal, etc.)
\usepackage{accvabbrv}

% Include other packages here, before hyperref.
\usepackage{graphicx}
\usepackage{booktabs}

% The "axessiblity" package can be found at: https://ctan.org/pkg/axessibility?lang=en
\usepackage[accsupp]{axessibility}  % comment out if not installed  % Improves PDF readability for those with disabilities.


% ---------------------------------------------------------------
% Hyperref package
% If you comment hyperref and then uncomment it, you should delete
% egpaper.aux before re-running latex.  (Or just hit 'q' on the first latex
% run, let it finish, and you should be clear).

\definecolor{accvblue}{rgb}{0.12,0.49,0.85}
\usepackage[pagebackref,breaklinks,colorlinks,citecolor=accvblue]{hyperref}


% ---------------------------------------------------------------
% If you wish to avoid re-using figure, table, and equation numbers from
% the main paper, please uncomment the following and change the numbers
% appropriately.
%\setcounter{figure}{2}
%\setcounter{table}{1}
%\setcounter{equation}{2}

% If you wish to avoid re-using reference numbers from the main paper,
% please uncomment the following and change the counter for `enumiv' to
% the number of references you have in the main paper (here, 6).
%\let\oldthebibliography=\thebibliography
%\let\oldendthebibliography=\endthebibliography
%\renewenvironment{thebibliography}[1]{%
%     \oldthebibliography{#1}%
%     \setcounter{enumiv}{6}%
%}{\oldendthebibliography}


% ---------------------------------------------------------------
% Paper ID for submission 230.
\def\paperID{230}
\def\confName{ACCV}
\def\confYear{2026}

\begin{document}

\title{Measuring and Mitigating Spatiotemporal Aliasing in Video Action
Recognition: A Cross-Architecture Analysis at Scale}

\maketitle
\thispagestyle{empty}
\appendix

We thank the reviewers, and we ran every experiment requested.

\paragraph{What is unchanged.}
\textsc{TDS}, our main contribution, is unaffected. \textbf{R2} asks whether its
ranking is an artefact of the endpoint definition; it is not. Integrating the
full stride curve rather than its endpoints gives an identical dataset ranking
(Spearman $1.000$); removing the positive part is \emph{numerically} identical,
since it never activates; normalising by baseline accuracy or by headroom above
chance exchanges only adjacent pairs ($0.952$). The architecture ordering also
holds: supplying every architecture the same distinct frames at the same
positions --- the matched comparison \textbf{R2} asks for --- preserves it at
Spearman $0.857$ ($p{=}0.007$, all 8 models). We withdraw only the
$3$--$5\times$ magnitude, which tracks input length rather than temporal
aggregation, and restate the finding as evidence efficiency: \emph{at an equal
budget of distinct frames, divided-attention and state-space designs deliver
more accuracy per frame}.

\paragraph{Characterising the temporal axis (R2, R3).}
\textbf{R3} is right that ``aliasing'' was imprecise for what we vary. Strided
candidates are re-uniformised to each model's fixed frame budget $B$, so at
fixed coverage the selected frames span the same window at every stride, and the
input changes only once $\lceil\text{window}/s\rceil<B$. \cref{tab:seconds}
shows this in seconds between sampled frames: four of five stride settings are
the same sampling rate on UCF-101. The axis therefore measures available
evidence, and we report it in seconds throughout, as \textbf{R3} suggests.

\begin{table}[t]
\centering\footnotesize
\setlength{\tabcolsep}{4pt}
\begin{tabular}{lccccc}
\toprule
Median s between frames & $s{=}1$ & $s{=}2$ & $s{=}4$ & $s{=}8$ & $s{=}16$ \\
\midrule
UCF-101 & 0.395 & 0.395 & 0.395 & 0.395 & 0.617 \\
AUTSL   & 0.125 & 0.125 & 0.131 & 0.252 & 0.475 \\
\bottomrule
\end{tabular}
\caption{Realised sampling interval at coverage 100\%, measured frame rates.}
\label{tab:seconds}
\end{table}

Two consequences follow. Pinning the frame count and varying only the window
width --- which changes the realised rate tenfold --- \emph{raises} accuracy in
7 of 8 datasets (Spearman $+1.0$ in four), so the binding constraint over this
range is temporal context, exactly \textbf{R3}'s reading. And the correlation
between input length and measured degradation, $0.849$ ($p{=}0.008$) under the
original protocol, falls to $0.231$ ($p{=}0.58$) once evidence is matched.

\paragraph{R1: multi-stride augmentation.}
We ran the ablation you proposed. \textbf{Your hypothesis is partly right, and
more than we expected}: exposure recovers $28.4\%$ of the drop for TimeSformer
and $37.0\%$ for R3D-18, whose stride-16 accuracy rises $6.5\%\!\to\!31.9\%$. It
is a trade rather than a repair --- 63--72\% persists and dense accuracy costs
$3.3$/$9.9$pp --- and at matched dense accuracy, training through the degraded
input recovers $2.4$--$6.5\times$ more than uniform resampling. On convergence:
across 730 training logs, 48px is indeed where the budget binds ($52.8\%$ of
runs peak at the final epoch \vs 16--27\% elsewhere), though the last three
epochs win only $1.38$pp; we state this as a limitation. Sec.~3.3 of the paper now gives the resize and PE-interpolation procedure, and Table~7 gains the
without-fine-tuning row.

\paragraph{R3: further items.}
Frame rates are now measured, not assumed --- three datasets differed (HMDB-51
30, AUTSL 30, EPIC 50) --- and at fixed nominal stride the realised interval
varies $3.9\times$ across datasets. Kinetics-400 is added, requiring no training
since all backbones are K400-pretrained; it sits at the low-demand end, which
sharpens the paper's point. Table~5 gains FineGym with the tier definition stated
explicitly, Table~7 the baseline column, and model ordering is now consistent
throughout. On Table~7 you are right that native-resolution tuning already gains
$+6$ to $+10.6$pp, so the $+39.2$pp separates into $\approx$32pp adaptation and
$\approx$7pp additional training; we report both.

\paragraph{R2: routing and statistics.}
Both points are correct. The rule thresholds $\max_c p(c|v)$, not entropy, and is
renamed. Under the held-out calibration you require (1000 stratified splits),
SSv2 moves from $+0.21$pp to $-0.03$pp (95\% CI $[-1.07,+0.83]$): the cascade
\emph{matches} dense accuracy at $6.9$ frames. We remove the corresponding
abstract claim and ``Mitigating'' from the title. Re-running the ANOVA as a
within-subject design gives partial $\eta^2$ of $0.292$ (coverage), $0.232$
(stride) and $0.070$ (interaction), so coverage leads by $1.26\times$ rather than
$2.2\times$ while the per-model ordering is unchanged; the interaction is now
given for all 56 pairs. The device-scaling factors were estimates and are
removed.

\end{document}
